\scanpage{132}
% archive foliation: 131
Quand on a de telles données algébriques, sur courbe algébrique $X/\mathbb C\inserted{\text{ sur }\mathbb C}$, comment reconstituer la structure analytique du th. en termes de celles-ci ? Il faut donc construire \inserted{l'ouvert} $\widetilde X\subset(P_X)^{\mathrm{an}}$ et $\widetilde X\longrightarrow(X_X)^{\mathrm{an}}$.

Or on a déjà un hom. formel d), et celui-ci pourra se prolonger, suivant les cas : $1^\circ)$ en $P_X\longrightarrow X_X$ par un \uline{iso} (cas $X$ droite projective). $2^\circ)$ Un \deleted{\unsure{germe}}\inserted{\unsure{hom.}}
\[
\deleted{(P_X^{\mathrm{an}}-\sigma'(X))}\quad
\inserted{\widetilde X^{\mathrm{alg}}}\unsure{=}
P_X^{\mathrm{an}}-\sigma'^{\mathrm{an}}(X^{\mathrm{an}})\xrightarrow{p}(X_X)^{\mathrm{an}},
\]
où $\sigma'$ est une section horizontale (\deleted{\unsure{unique}} toujours disjointe de $\sigma$) de $P_X$ sur $X$, faisant de $\widetilde X$ un rev. de $(X_X)^{\mathrm{an}}$, et transformant $\sigma$ en $\Delta_X$. (\deleted{\unsure{Mais}} $p$ est un iso ssi $X\simeq \mathbb E^1_{\mathbb C}$ i.e. $X$ simplement connexe non propre). Ce cas se présente pour $X\simeq \mathbb E^1_{\mathbb C}$, $\mathbb G_{m\mathbb C}$, ou $X$ courbe elliptique.

\margin{\deleted{NB \unsure{Par Deligne le faisceau des sections est algébrique !}}}

\uline{Question} A-t-on une description \uline{purement algébrique} de $\sigma'$ (\unsure{en supposant qu'on se soit muni de} $(P_X,c_X,\unsure{\rho_X|\sigma_X}))$ ? Et de $p_X$ dans le cas où c'est un iso analytique i.e. $X\simeq \mathbb E^1_{\mathbb C}$ --- alors $p_X$ est lui-même algébrique !

\margin{\uline{OK} on trouve \unsure{en $2^\circ$ un groupe vectoriel} $G$, et \unsure{en prenant la clôture alg.} \unsure{pour avoir la clôture projective du fibré} \unsure{et des sections horizontales}.}

$3^\circ)$ On a un ouvert $\widetilde X\deleted{\unsure{(qui est un rev.)}}\inserted{\text{dont les fibres sont des}}$ disques des fibres de $P_X^{\mathrm{an}}$, lesquels disques sont les domaines d'existence maximaux des
\[\widehat p_x:(P_X)_{x\mathrm{an}}\longrightarrow X^{\mathrm{an}}\ldots\]
